\documentclass[11pt,a4paper]{ltjsarticle}
\usepackage{amsmath,amssymb}
\usepackage{graphicx}
\usepackage{booktabs}
\usepackage{array}
\usepackage{luatexja-fontspec}
\setmainjfont{HaranoAjiMincho-Regular}
\setsansjfont{HaranoAjiGothic-Medium}
\linespread{1.05}\selectfont

% Recommended PDF filename:
% アルゴリズム設計論_repo5_2611193_中井平蔵.pdf
% Example build command:
% latexmk -lualatex -jobname=アルゴリズム設計論_repo5_2611193_中井平蔵 main.tex

\begin{document}

\begin{center}
{\LARGE \textbf{アルゴリズム設計論 第5回レポート}}
\end{center}
\vspace{1em}

\begin{flushright}
提出日: \today \\
学籍番号: 2611193 \\
氏名: 中井 平蔵
\end{flushright}

\medskip
問題文より、与えられた品物の重さと価値は以下の表であり、ナップサックの容量は \(15\) である。

\[
\renewcommand{\arraystretch}{1.2}
\begin{array}{c|cccccc}
\hline
\text{品物番号} & 1 & 2 & 3 & 4 & 5 & 6 \\
\hline
\text{重さ} & 8 & 3 & 6 & 1 & 10 & 5 \\
\text{価値} & 1 & 4 & 2 & 3 & 1 & 7 \\
\hline
\end{array}
\]

F(h, g) を、品物 \(1\) から品物 \(h\) までを使用し、容量 \(g\) 以下で得られる価値の最大値と定義する。品物 \(h\) の重さを \(w_h\)、価値を \(v_h\) とする。

\begin{itemize}
\item 品物 \(h\) の重さが容量 \(g\) を超える場合：
\[
F(h, g) = F(h - 1, g)
\]
\item 品物 \(h\) を入れることができる場合：
\[
F(h, g) = \max\bigl(v_h + F(h - 1, g - w_h),\; F(h - 1, g)\bigr)
\]
\end{itemize}

これを動的計画法の表にまとめると、次のようになる。

\[
\renewcommand{\arraystretch}{1.2}
\begin{array}{c|cccccccccccccccc}
\hline
\shortstack[c]{\(g\)} & 0 & 1 & 2 & 3 & 4 & 5 & 6 & 7 & 8 & 9 & 10 & 11 & 12 & 13 & 14 & 15 \\
\hline
F(0,g) & 0 & 0 & 0 & 0 & 0 & 0 & 0 & 0 & 0 & 0 & 0 & 0 & 0 & 0 & 0 & 0 \\
F(1,g) & 0 & 0 & 0 & 0 & 0 & 0 & 0 & 0 & 1 & 1 & 1 & 1 & 1 & 1 & 1 & 1 \\
F(2,g) & 0 & 0 & 0 & 4 & 4 & 4 & 4 & 4 & 4 & 4 & 4 & 5 & 5 & 5 & 5 & 5 \\
F(3,g) & 0 & 0 & 0 & 4 & 4 & 4 & 4 & 4 & 4 & 6 & 6 & 6 & 6 & 6 & 6 & 6 \\
F(4,g) & 0 & 3 & 3 & 4 & 7 & 7 & 7 & 7 & 7 & 7 & 9 & 9 & 9 & 9 & 9 & 9 \\
F(5,g) & 0 & 3 & 3 & 4 & 7 & 7 & 7 & 7 & 7 & 7 & 9 & 9 & 9 & 9 & 9 & 9 \\
F(6,g) & 0 & 3 & 3 & 4 & 7 & 7 & 10 & 10 & 11 & 14 & 14 & 14 & 14 & 14 & 14 & 16 \\
\hline
\end{array}
\]

表より、\(F(6, 15) = 16\) であるため、最大価値は \(16\) である。

表の右下 \(F(6, 15)\) から逆にたどることで、価値が最大となる品物の組み合わせを確認できる。

\begin{enumerate}
\item \(F(6, 15) = 16\)、\(F(5, 15) = 9\) なので、品物6を選ぶ。

残り容量は
\[
15 - 5 = 10
\]

\item \(F(5, 10) = F(4, 10) = 9\) なので、品物5は選ばない。

\item \(F(4, 10) = 9\)、\(F(3, 10) = 6\) なので、品物4を選ぶ。

残り容量は
\[
10 - 1 = 9
\]

\item \(F(3, 9) = 6\)、\(F(2, 9) = 4\) なので、品物3を選ぶ。

残り容量は
\[
9 - 6 = 3
\]

\item \(F(2, 3) = 4\)、\(F(1, 3) = 0\) なので、品物2を選ぶ。

残り容量は
\[
3 - 3 = 0
\]
\end{enumerate}

したがって、選択した品物は \fbox{\(\{2, 3, 4, 6\}\)}、合計重量は
\(3 + 6 + 1 + 5 = \fbox{\(15\)}\)、合計価値は \fbox{\(16\)} である。

\end{document}
